\documentclass[aspectratio=169]{beamer}

% Essential packages (core)
\usepackage{hyperref}

    % Configure hyperref for better citation links
    \hypersetup{
        colorlinks=true,
        linkcolor=blue,
        citecolor=blue,
        urlcolor=blue,
        linktoc=all,
        pdfborder={0 0 0}
    }
    
\usepackage{graphicx}
\usepackage{amsmath}
\usepackage{tikz}
\usepackage{pgfplots}
\usepackage{xstring}
\usepackage{animate}
\usepackage{multimedia}
\usepackage{xifthen}
\usepackage{xcolor}
\usepackage{geometry}
\usepackage{booktabs}
\usepackage{grffile}
\geometry{paperwidth=128mm,paperheight=96mm}

% Load TikZ libraries FIRST
\usetikzlibrary{positioning}
\usetikzlibrary{shapes.symbols}
\usetikzlibrary{shapes.callouts}
\usetikzlibrary{shapes.multipart}
\usetikzlibrary{calc}
\usetikzlibrary{overlay-beamer-styles}
\usetikzlibrary{shapes.geometric}
\usetikzlibrary{arrows.meta}
\usetikzlibrary{backgrounds}
\usetikzlibrary{fit}

% Define the style for covered text
\setbeamercovered{dynamic}
\setbeamerfont{item projected}{size=\small}
\setbeamercolor{alerted text}{fg=white}

% Extended packages with fallbacks
\IfFileExists{tcolorbox.sty}{\usepackage{tcolorbox}}{}
\IfFileExists{fontawesome5.sty}{\usepackage{fontawesome5}}{}
\IfFileExists{pifont.sty}{\usepackage{pifont}}{}
\IfFileExists{soul.sty}{\usepackage{soul}}{}

% Package configurations
\pgfplotsset{compat=1.18}

% SIMPLIFIED text effects - use standard LaTeX instead of problematic TikZ
\newcommand{\shadowtext}[2][2pt]{%
    \textcolor{white}{\textbf{#2}}%
}

\newcommand{\glowtext}[2][myblue]{%
    \textcolor{#1}{\textbf{#2}}%
}

% Conditional definitions based on package availability
% CRITICAL: Use ##1 instead of #1 inside \IfFileExists
\IfFileExists{tcolorbox.sty}{%
    \newtcolorbox{alertbox}[1][red]{%
        colback=##1!5!white,
        colframe=##1!75!black,
        fonttitle=\bfseries,
        boxrule=0.5pt,
        rounded corners
    }

    \newtcolorbox{infobox}[1][blue]{%
        enhanced,
        colback=##1!5!white,
        colframe=##1!75!black,
        arc=4mm,
        boxrule=0.5pt,
        fonttitle=\bfseries,
        attach boxed title to top center={yshift=-3mm,yshifttext=-1mm},
        boxed title style={size=small,colback=##1!75!black}
    }
}{}

% Define colors
\definecolor{myred}{RGB}{255,50,50}
\definecolor{myblue}{RGB}{0,130,255}
\definecolor{mygreen}{RGB}{0,200,100}
\definecolor{myyellow}{RGB}{255,210,0}
\definecolor{myorange}{RGB}{255,130,0}
\definecolor{mypurple}{RGB}{147,112,219}
\definecolor{mypink}{RGB}{255,105,180}
\definecolor{myteal}{RGB}{0,128,128}
\definecolor{mygray}{RGB}{128,128,128}
\definecolor{mybrown}{RGB}{139,69,19}
\definecolor{mycyan}{RGB}{0,255,255}

% Basic highlighting commands
\newcommand{\hlbias}[1]{\textcolor{myblue}{\textbf{#1}}}
\newcommand{\hlvariance}[1]{\textcolor{mypink}{\textbf{#1}}}
\newcommand{\hltotal}[1]{\textcolor{myyellow}{\textbf{#1}}}
\newcommand{\hlkey}[1]{\colorbox{myblue!20}{\textbf{#1}}}
\newcommand{\hlnote}[1]{\colorbox{mygreen!20}{\textbf{#1}}}

% ========== LAYOUT COMMANDS - Check if already defined ==========
% Split layout: image on left, text on right
\ifcsname split\endcsname\else
\newcommand{\split}[2]{%
    \begin{columns}[T]
        \begin{column}{0.45\textwidth}
            \begin{center}
                \includegraphics[width=\textwidth,keepaspectratio]{#1}
            \end{center}
        \end{column}
        \begin{column}{0.5\textwidth}
            #2
        \end{column}
    \end{columns}
}
\fi

% PIP layout: text on left, small image on right
\ifcsname pip\endcsname\else
\newcommand{\pip}[2]{%
    \begin{columns}[T]
        \begin{column}{0.68\textwidth}
            #2
        \end{column}
        \begin{column}{0.28\textwidth}
            \vspace{1em}
            \includegraphics[width=\textwidth,keepaspectratio]{#1}
        \end{column}
    \end{columns}
}
\fi

% Fullframe layout: image takes entire frame
\ifcsname ff\endcsname\else
\newcommand{\ff}[1]{%
    \setbeamertemplate{background}{%
        \begin{tikzpicture}[remember picture,overlay]
            \node at (current page.center) {%
                \includegraphics[width=\paperwidth,height=\paperheight,keepaspectratio]{#1}
            };
        \end{tikzpicture}%
    }
    \begin{center}
        \vfill
        \textcolor{white}{\textbf{Full Frame Image}}
        \vfill
    \end{center}
}
\fi

% Watermark layout: image as background watermark
\ifcsname wm\endcsname\else
\newcommand{\wm}[1]{%
    \setbeamertemplate{background}{%
        \begin{tikzpicture}[remember picture,overlay]
            \node[opacity=0.15] at (current page.center) {%
                \includegraphics[width=\paperwidth,height=\paperheight,keepaspectratio]{#1}
            };
        \end{tikzpicture}%
    }
}
\fi

% Highlight layout: image with highlighted content overlay
\ifcsname hl\endcsname\else
\newcommand{\hl}[2]{%
    \begin{columns}[T]
        \begin{column}{0.6\textwidth}
            \includegraphics[width=\textwidth,keepaspectratio]{#1}
        \end{column}
        \begin{column}{0.36\textwidth}
            \colorbox{yellow!20}{%
                \begin{minipage}{\textwidth}
                    #2
                \end{minipage}%
            }
        \end{column}
    \end{columns}
}
\fi

% Background layout: image as background with content overlay
\ifcsname bg\endcsname\else
\newcommand{\bg}[1]{%
    \setbeamertemplate{background}{%
        \begin{tikzpicture}[remember picture,overlay]
            \node[opacity=0.3] at (current page.center) {%
                \includegraphics[width=\paperwidth,height=\paperheight,keepaspectratio]{#1}
            };
        \end{tikzpicture}%
    }
}
\fi

% Top-bottom layout: image at top, content at bottom
\ifcsname tb\endcsname\else
\newcommand{\tb}[2]{%
    \begin{center}
        \includegraphics[width=0.8\textwidth,keepaspectratio]{#1}
    \end{center}
    \vspace{0.5em}
    #2
}
\fi

% Overlay layout: image with overlaid text blocks
\ifcsname ol\endcsname\else
\newcommand{\ol}[1]{%
    \begin{tikzpicture}[remember picture,overlay]
        \node at (current page.center) {%
            \includegraphics[width=\paperwidth,keepaspectratio]{#1}
        };
    \end{tikzpicture}%
}
\fi

% Corner layout: image in corner
\ifcsname corner\endcsname\else
\newcommand{\corner}[2]{%
    \begin{tikzpicture}[remember picture,overlay]
        \node[anchor=south east] at (current page.south east) {%
            \includegraphics[width=0.25\textwidth,keepaspectratio]{#1}
        };
    \end{tikzpicture}%
    #2
}
\fi

% Mosaic layout: grid of images
\ifcsname mosaic\endcsname\else
\newcommand{\mosaic}[2]{%
    \begingroup
    \def\mosaic@params{#1}%
    \def\mosaic@images{#2}%
    \pgfmathsetmacro{\mosaic@rows}{{\mosaic@params}[0]}%
    \pgfmathsetmacro{\mosaic@cols}{{\mosaic@params}[2]}%
    \begin{center}
    \begin{tabular}{*{\mosaic@cols}{c}}
    \hline
    \mosaic@process
    \hline
    \end{tabular}
    \end{center}
    \endgroup
}

\def\mosaic@process{%
    \mosaic@process@helper\mosaic@images,\@empty
}

\def\mosaic@process@helper#1,#2\@empty{%
    \ifx\@empty#2\@empty
        \includegraphics[width=0.3\textwidth,keepaspectratio]{#1}%
    \else
        \includegraphics[width=0.3\textwidth,keepaspectratio]{#1} &
        \def\mosaic@remaining{#2}%
        \mosaic@process@next
    \fi
}

\def\mosaic@process@next{%
    \mosaic@process@helper\mosaic@remaining,\@empty
}
\fi

% Basic theme setup
\usetheme{Madrid}
\usecolortheme{owl}

% Color settings
\setbeamercolor{normal text}{fg=white}
\setbeamercolor{structure}{fg=myyellow}
\setbeamercolor{alerted text}{fg=myorange}
\setbeamercolor{example text}{fg=mygreen}
\setbeamercolor{background canvas}{bg=black}
\setbeamercolor{frametitle}{fg=white,bg=black}

% Notes support
\usepackage{pgfpages}
\setbeameroption{show notes on second screen=right}
\setbeamertemplate{note page}{\pagecolor{yellow!5}\insertnote}

% Animated background support
\newcommand{\anbg}[2][0.2]{%
    \ifx\@empty#2\@empty
        % Clear background if empty argument
        \setbeamertemplate{background}{}
    \else
        % Set animated background
        \setbeamertemplate{background}{%
            \begin{tikzpicture}[remember picture,overlay]
                \node[opacity=#1] at (current page.center) {%
                    \animategraphics[autoplay,loop,width=\paperwidth]{12}{#2}{}{}
                };
            \end{tikzpicture}%
        }
    \fi
}


% Progress bar setup
\makeatletter
\def\progressbar@progressbar{}
\newcount\progressbar@tmpcounta
\newcount\progressbar@tmpcountb
\newdimen\progressbar@pbht
\newdimen\progressbar@pbwd
\newdimen\progressbar@tmpdim

\progressbar@pbwd=\paperwidth
\progressbar@pbht=2pt

\def\progressbar@progressbar{%
   \begin{tikzpicture}[very thin]
       \ifnum\insertframenumber>0
           \pgfmathparse{\insertframenumber/\inserttotalframenumber}
           \edef\progress@ratio{\pgfmathresult}
           \shade[top color=myblue!50,bottom color=myblue]
               (0pt, 0pt) rectangle (\progress@ratio\progressbar@pbwd, \progressbar@pbht);
       \fi
   \end{tikzpicture}%
}

% Modified frame title template
\setbeamertemplate{frametitle}{
   \nointerlineskip
   \vskip1ex
   \begin{beamercolorbox}[wd=\paperwidth,ht=4ex,dp=2ex]{frametitle}
       \begin{minipage}[t]{\dimexpr\paperwidth-4em}
           \centering
           \vspace{2pt}
           \insertframetitle
           \vspace{2pt}
       \end{minipage}
   \end{beamercolorbox}
   \vskip.5ex
   \progressbar@progressbar
}
\makeatother
\makeatletter\def\insertshortinstitute{airis4D}\makeatother

% Footline template
\setbeamertemplate{footline}{%
 \leavevmode%
 \hbox{%
   \begin{beamercolorbox}[wd=.333333\paperwidth,ht=2.25ex,dp=1ex,center]{author in head/foot}%
     \usebeamerfont{author in head/foot}\insertshortauthor~(\insertshortinstitute)%
   \end{beamercolorbox}%
   \begin{beamercolorbox}[wd=.333333\paperwidth,ht=2.25ex,dp=1ex,center]{title in head/foot}%
     \usebeamerfont{title in head/foot}\insertshorttitle%
   \end{beamercolorbox}%
   \begin{beamercolorbox}[wd=.333333\paperwidth,ht=2.25ex,dp=1ex,right]{date in head/foot}%
     \usebeamerfont{date in head/foot}\insertshortdate{}\hspace*{2em}%
     \insertframenumber{} / \inserttotalframenumber\hspace*{2ex}%
   \end{beamercolorbox}}%
 \vskip0pt%
}

% Additional settings
\setbeamersize{text margin left=5pt,text margin right=5pt}
\setbeamertemplate{navigation symbols}{}
\setbeamertemplate{blocks}[rounded][shadow=true]
% Title setup
\title{The Sloan Digital Sky Survey}
\subtitle{SDSS}
\author{Ninan Sajeeth Philip}
\institute{\textcolor{mygreen}{Artificial Intelligence Research and Intelligent Systems (airis4D)}}
\date{\today}

    % Additional packages for special characters and units
    \usepackage{textcomp}      % For \textmu, \textendash, etc.
    \usepackage{siunitx}       % For proper units (Ω·cm², etc.)
    \usepackage{amsmath}       % Enhanced math support
    \usepackage{amssymb}       % Additional math symbols

    % Configure siunitx for proper formatting
    \sisetup{
        per-mode = symbol,
        output-decimal-marker = {.},
        group-separator = {,}
    }


    % Color Information for TikZ Figures
    % Use these color definitions for consistent theming
    % Text colors are automatically chosen using XOR rule

    % Default color palette
    \definecolor{airis4d_blue}{RGB}{41,128,185}
    \definecolor{airis4d_green}{RGB}{39,174,96}
    \definecolor{airis4d_orange}{RGB}{243,156,18}
    \definecolor{airis4d_red}{RGB}{231,76,60}
    \definecolor{airis4d_purple}{RGB}{155,89,182}
    \definecolor{airis4d_teal}{RGB}{26,188,156}
    \definecolor{airis4d_gray}{RGB}{149,165,166}

    % XOR Rule for text colors:
    % - Use text=white on dark backgrounds (luminance < 0.179)
    % - Use text=black on light backgrounds (luminance > 0.179)

    % Example usage:
    % \node[fill=airis4d_blue, text=white] {White text on blue};
    % \node[fill=airis4d_orange, text=black] {Black text on orange};
    
\begin{document}

\begin{frame}[plain]
\titlepage
\end{frame}

\begin{frame}{Abstract \& Objectives}
\frametitle{Abstract \& Objectives}

\textbf{Abstract \& Objectives}
\begin{block}{The "Artificial Photosynthesis" Model}
DSSCs mimic photosynthetic systems: a molecular dye harvests light and injects electrons into a wide-bandgap semiconductor, separating charge without a traditional p-n junction.
\end{block}
\textbf{Objectives:}
\begin{itemize}
\item Understand the photoelectrochemical operating principle.
\item Analyze energy-level alignment, kinetics, and recombination.
\item Review materials: sensitizers, redox mediators, electrodes.
\item Discuss efficiency limits, stability, and emerging applications.
\end{itemize}
\vspace{2mm}
\cite{Grätzel2001, Hagfeldt2010}
\end{frame}


\begin{frame}{Evolution of Solar Technology}
\frametitle{Evolution of Solar Technology}

\textbf{Evolution of Solar Technology} — 1st to 3rd Generation
\begin{table}
\centering
\begin{tabular}{lll}
\toprule
Generation & Technology & Key Limitation \\
\midrule
1st & Crystalline Si (mono/poly) & High cost, rigid \\
2nd & Thin-film (CIGS, CdTe) & Scarce/toxic materials \\
3rd & DSSC, Perovskites, OPV & Stability (DSSC: electrolyte leakage) \\
\bottomrule
\end{tabular}
\end{table}
\begin{itemize}
\item DSSCs: Low-cost, flexible, work under diffuse light.
\item Record lab efficiency (DSSC): $\sim$13\% (2022, co-sensitized).
\end{itemize}
\cite{Green2022}
\end{frame}


\begin{frame}{The Breakthrough (1991)}
\frametitle{The Breakthrough (1991)}

\textbf{The Breakthrough (1991)} — O'Regan \& Grätzel, \textit{Nature
}
\begin{quote}
"A low-cost, high-efficiency solar cell based on dye-sensitized colloidal $TiO_2$ films"
\end{quote}
\textbf{Key innovations:}
\begin{itemize}
\item Mesoporous $TiO_2$ (10–20 nm particles) → high surface area.
\item Ruthenium polypyridyl dye (efficient visible absorption).
\item Iodine/triiodide redox electrolyte.
\item Achieved 7.1–7.9\% efficiency under AM 1.5.
\end{itemize}
\textbf{Impact:} Launched third-generation photovoltaics; $>$30,000 citations.
\cite{ORegan1991}
\end{frame}


\begin{frame}{Defining DSSCs}
\frametitle{Defining DSSCs}

\textbf{Defining DSSCs} — Photoelectrochemical vs. Solid-State
\begin{columns}
\column{0.5\textwidth}
\textbf{DSSC = Photoelectrochemical system}
\begin{itemize}
\item Charge separation at dye/electrolyte interface.
\item No depletion layer; diffusion-driven transport.
\item Liquid or solid-state redox mediator.
\end{itemize}
\column{0.5\textwidth}
\textbf{Conventional p-n junction}
\begin{itemize}
\item Built-in electric field.
\item Drift-dominated charge extraction.
\item Rigid architecture.
\end{itemize}
\end{columns}
\vspace{3mm}
\textit{Result:} DSSCs tolerate defects, work indoors, and offer transparency.
\cite{Peter2007}
\end{frame}


\begin{frame}{General Anatomy}
\frametitle{General Anatomy}

\textbf{General Anatomy} — The "Sandwich" Structure
\begin{center}
\fbox{\parbox{0.85\textwidth}{\centering
Glass / FTO substrate \\
$\uparrow$ \\
Mesoporous TiO$_2$ + Dye monolayer (Photoanode) \\
$\uparrow$ \\
Redox Electrolyte ($I^-/I_3^-$ or Co$^{\text{II/III}}$) \\
$\uparrow$ \\
Platinum / Carbon (Counter Electrode) \\
$\uparrow$ \\
Glass / FTO substrate
}}
\end{center}
Thickness: $\sim$10 \textmu m TiO$_2$, $\sim$2 mm electrolyte gap (lab cells).
\cite{Kalyanasundaram2010}
\end{frame}


\begin{frame}{The Substrate}
\frametitle{The Substrate}

\textbf{The Substrate} — FTO and ITO Requirements
\begin{itemize}
\item \textbf{FTO (Fluorine-doped Tin Oxide):}
\begin{itemize}
\item High thermal stability (up to 500°C for sintering).
\item Sheet resistance $\sim$ 7–15 $\Omega$/sq.
\item Preferred for DSSCs.
\end{itemize}
\begin{itemize}
\item \textbf{ITO (Indium Tin Oxide):}
\end{itemize}
\begin{itemize}
\item Higher conductivity but degrades under sintering.
\item Used for flexible, low-temperature DSSCs.
\end{itemize}
\end{itemize}
Transmittance $>$80\% in visible needed for front illumination.
\cite{Gordillo2006}
\end{frame}


\begin{frame}{The Photoanode}
\frametitle{The Photoanode}

\textbf{The Photoanode} — Mesoporous Metal Oxides
\textbf{Common oxides:}
\begin{itemize}
\item \textbf{TiO$_2$ (anatase):} Bandgap $\sim$3.2 eV; CB edge $\sim$ -0.5 V vs. NHE.
\item \textbf{ZnO:} Higher electron mobility but lower chemical stability.
\item \textbf{SnO$_2$:} Higher CB edge → lower $V_{oc}$; used in some high-voltage designs.
\end{itemize}
\textbf{Key property:} Mesoporosity (pore size 15–30 nm) allows electrolyte penetration and dye loading.
\cite{Hagfeldt2010}
\end{frame}


\begin{frame}{Role of Nanostructure}
\frametitle{Role of Nanostructure}

\textbf{Role of Nanostructure} — High Surface-to-Volume Ratio
\begin{itemize}
\item Planar TiO$_2$: Dye loading $\sim$ 0.1 nmol/cm$^2$ → negligible current.
\item Mesoporous film (10 \textmu m thick, 20 nm particles): Roughness factor $\sim$ 1000 → loading $\sim$ 100 nmol/cm$^2$.
\end{itemize}
\textbf{Consequence:} Optical density $\gg$1 possible even with monolayer dye coverage.
\vspace{2mm}
\textit{Why not thicker?} Increased recombination and mass transport limitations.
\cite{Grätzel2005}
\end{frame}


\begin{frame}{The Sensitizer (Dye)}
\frametitle{The Sensitizer (Dye)}

\textbf{The Sensitizer (Dye)} — Molecular Interface
\textbf{Functions:}
\begin{enumerate}
\item Absorb photons (350–900 nm ideally).
\item Inject electrons into TiO$_2$ CB ($\tau_{inj} \sim$ fs–ps).
\item Be regenerated by redox mediator before recombination.
\end{enumerate}
\textbf{Types:}
\begin{itemize}
\item Ru-complexes (N719, Black dye) – high efficiency, expensive.
\item Organic D-$\pi$-A – high extinction, tunable.
\item Porphyrins – strong absorbers, used in record cells.
\end{itemize}
\cite{Mathew2014}
\end{frame}


\begin{frame}{The Electrolyte}
\frametitle{The Electrolyte}

\textbf{The Electrolyte} — Redox Mediator
\textbf{Classic system:} Iodine/triiodide ($I^-/I_3^-$) in acetonitrile.
\begin{itemize}
\item $E^0 \approx +0.35$ V vs. NHE.
\item Fast dye regeneration ($\mu$s).
\item Limitation: Corrodes silver contacts, absorbs blue light.
\end{itemize}
\textbf{Alternatives:}
\begin{itemize}
\item Cobalt complexes ($\text{Co}^{\text{II/III}}$): Higher $V_{oc}$ but slower regeneration.
\item Copper complexes ($\text{Cu}^{\text{I/II}}$): Emerging high-performance.
\end{itemize}
\cite{Feldt2010}
\end{frame}


\begin{frame}{Counter Electrode}
\frametitle{Counter Electrode}

\textbf{Counter Electrode} — Catalyst Requirements
\textbf{Role:} Reduce $I_3^-$ to $I^-$ (or equivalent) efficiently.
\begin{itemize}
\item \textbf{Platinum:} Best catalytic activity; high cost.
\item \textbf{Carbon materials:} Carbon black, graphene, CNTs – stable, cheap.
\item \textbf{Conducting polymers:} PEDOT:PSS – flexible, transparent.
\end{itemize}
\textit{Requirements:} Low charge-transfer resistance ($R_{ct} < 5\ \Omega\cdot$cm$^2$), good adhesion.
\cite{Trancik2007}
\end{frame}


\begin{frame}{The 5-Step Operational Cycle}
\frametitle{The 5-Step Operational Cycle}

\textbf{The 5-Step Operational Cycle} — From Absorption to Regeneration
\begin{enumerate}
\item \textbf{Photoexcitation:} $D + h\nu \rightarrow D^*$
\item \textbf{Electron injection:} $D^* \rightarrow D^+ + e^- (\text{TiO}_2)$
\item \textbf{Transport:} $e^-$ diffuses through TiO$_2$ network to FTO.
\item \textbf{Dye regeneration:} $2D^+ + 3I^- \rightarrow 2D + I_3^-$
\item \textbf{Counter electrode reaction:} $I_3^- + 2e^- (\text{Pt}) \rightarrow 3I^-$
\end{enumerate}
Overall: $h\nu + e^- (\text{Pt}) \rightarrow e^- (\text{TiO}_2)$ no net chemical change.
\end{frame}


\begin{frame}{Step 1: Photoexcitation}
\frametitle{Step 1: Photoexcitation}

\textbf{Step 1: Photoexcitation} — $D + h\nu \rightarrow D^*$
\begin{itemize}
\item Absorbed photon promotes electron from HOMO to LUMO.
\item Singlet excited state ($^1D^*$) with lifetime 0.1–10 ns (organic) or 10–100 ns (Ru).
\item Must overlap solar spectrum: AM1.5G peak at $\sim$550 nm.
\end{itemize}
\textbf{Key metric:} Molar extinction coefficient $\varepsilon$ (M$^{-1}$cm$^{-1}$).
\begin{itemize}
\item Ru-complexes: $\varepsilon \approx 10^4$ at 530 nm.
\item Organic dyes: $\varepsilon \approx 2–5 \times 10^4$ – thinner films possible.
\end{itemize}
\cite{Ardo2009}
\end{frame}


\begin{frame}{Step 2: Electron Injection}
\frametitle{Step 2: Electron Injection}

\textbf{Step 2: Electron Injection} $\rightarrow D^* \rightarrow D^+ + e^{-CB}$
\textbf{Kinetics:}
\begin{itemize}
\item Ultrafast: typical $\tau_{inj} = 10$ fs – 100 ps.
\item Requires LUMO above TiO$_2$ CB edge (overpotential $\Delta E > 0.2$ eV).
\end{itemize}
\textbf{Mechanisms:}
\begin{itemize}
\item Adiabatic injection via anchoring group (carboxylate).
\item Non-adiabatic electron tunneling.
\end{itemize}
\textbf{Competing process:} Excited-state decay to ground state (fluorescence/heat).
\cite{Koops2009}
\end{frame}


\begin{frame}{Step 3: Charge Transport}
\frametitle{Step 3: Charge Transport}

\textbf{Step 3: Charge Transport} — Diffusion-Driven Network
\begin{itemize}
\item Electrons move by \textbf{multiple trapping} and release in TiO$_2$ nanoparticle network.
\item Effective diffusion coefficient $D_n \approx 10^{-4} – 10^{-3}$ cm$^2$/s (single crystal $D \approx 1$ cm$^2$/s).
\item Transport limited by trap states (surface oxygen vacancies).
\end{itemize}
\textbf{Consequence:} Films $>$20 \textmu m cause collection losses.
\vspace{2mm}
\emph{Model:} Random walk among particles; grain boundaries dominate.
\cite{Nelson2003}
\end{frame}


\begin{frame}{Step 4: Dye Regeneration}
\frametitle{Step 4: Dye Regeneration}

\textbf{Step 4: Dye Regeneration} — $2D^+ + 3I^- \rightarrow 2D + I_3^-$
\begin{itemize}
\item Oxidized dye ($D^+$) must be reduced before recombination with TiO$_2$ electrons.
\item Typical regeneration time: $\tau_{reg} \approx 1–10\ \mu$s (with $I^-$).
\item Overpotential for regeneration: $E_{\text{HOMO}}(D) - E^0(I^-/I_3^-) \approx 0.4–0.6$ eV.
\end{itemize}
\textbf{Failure mode:} If $\tau_{reg} > \tau_{rec}$ (recombination), efficiency drops.
\cite{Clifford2011}
\end{frame}


\begin{frame}{Step 5: Electrolyte Regeneration at Counter Electrode}
\frametitle{Step 5: Electrolyte Regeneration at Counter Electrode}

\textbf{Step 5: Electrolyte Regeneration at Counter Electrode} — $I_3^- + 2e^- \rightarrow 3I^-$
\begin{itemize}
\item Two-electron reduction occurs via Volmer–Heyrovský mechanism on Pt.
\item Fast kinetics: exchange current density $\sim$ mA/cm$^2$.
\item Poor catalysts (e.g., bare FTO) cause large overpotential → reduced $FF$.
\end{itemize}
\textbf{Design:} Sputtered Pt ($\sim$10 nm) or thermally decomposed H$_2$PtCl$_6$.
\cite{Hilgendorff2000}
\end{frame}


\begin{frame}{Kinetics}
\frametitle{Kinetics}

\textbf{Kinetics} — Time Scales of Forward vs. Recombination
\begin{table}
\centering
\begin{tabular}{lcc}
\toprule
Process & Time scale & Rate constant \\
\midrule
Injection & 10 fs – 100 ps & $k_{inj} \sim 10^{10} – 10^{13}$ s$^{-1}$ \\
Dye regeneration & 1 – 10 $\mu$s & $k_{reg} \sim 10^{5}$ s$^{-1}$ \\
Recombination (TiO$_2$ $e^-$ + $D^+$) & 10 $\mu$s – 1 ms & $k_{rec} \sim 10^{3}$ s$^{-1}$ \\
Recombination (TiO$_2$ $e^-$ + $I_3^-$) & 1 – 100 ms & $k_{ct} \sim 10$ s$^{-1}$ \\
\bottomrule
\end{tabular}
\end{table}
\textit{Golden rule:} $k_{inj} \gg k_{reg} \gg k_{rec}$ for high efficiency.
\cite{Durrant2004}
\end{frame}


\begin{frame}{Charge Extraction vs. Recombination}
\frametitle{Charge Extraction vs. Recombination}

\textbf{Charge Extraction vs. Recombination} — Diffusion Length Concept
\[
L_n = \sqrt{D_n \cdot \tau_n}
\]
\begin{itemize}
\item $D_n$: electron diffusion coefficient in TiO$_2$.
\item $\tau_n$: electron lifetime (determined by recombination with $I_3^-$).
\end{itemize}
\textbf{Criterion for efficient collection:} $L_n >$ film thickness ($\sim$10 \textmu m).
\vspace{2mm}
In good DSSCs: $D_n \approx 10^{-4}$ cm$^2$/s, $\tau_n \approx 10$ ms → $L_n \approx 30\ \mu$m.
\vspace{2mm}
Recombination reduces $\tau_n$ → shorter $L_n$ → lower $J_{sc}$.
\cite{Frank2004}
\end{frame}


\begin{frame}{Energy Level Alignment}
\frametitle{Energy Level Alignment}

\textbf{Energy Level Alignment} — LUMO, HOMO, Fermi Levels
\begin{center}
\fbox{\parbox{0.9\textwidth}{\centering
\vspace{1cm}
[Energy diagram: Vacuum level → CB TiO$_2$ (-0.5 V) → LUMO dye → HOMO dye → Redox potential (0.35 V) → VB TiO$_2$]
\vspace{1cm}
}}
\end{center}
$V_{oc} = E_{F,\text{TiO}_2} - E_{\text{redox}}$ (under illumination).
\cite{Cahen2003}
\end{frame}


\begin{frame}{Thermodynamics of Electron Injection}
\frametitle{Thermodynamics of Electron Injection}

\textbf{Thermodynamics of Electron Injection} — Overpotential Requirements
\begin{itemize}
\item Driving force: $\Delta G_{inj} = E_{\text{LUMO}}(\text{dye}) - E_{\text{CB}}(\text{TiO}_2)$ (must be negative).
\item Typically 0.2–0.4 eV overpotential for efficient injection.
\item Too large overpotential ($>$0.6 eV) wastes $V_{oc}$.
\end{itemize}
\textbf{Optimization:} Tune dye LUMO by electron-withdrawing anchoring group.
\vspace{2mm}
Example: N719 dye LUMO $\approx$ -1.1 V vs. NHE, CB TiO$_2$ $\approx$ -0.5 V → $\Delta G \approx -0.6$ eV.
\cite{Nazeeruddin2003}
\end{frame}


\begin{frame}{Open-Circuit Voltage ($V_{oc}$)}
\frametitle{Open-Circuit Voltage ($V_{oc}$)}

{Fermi Level of TiO$_2$ vs. Redox Potential}
At open circuit: $V_{oc} = \frac{kT}{q} \ln\left(\frac{J_{sc}}{J_0} + 1\right)$ but more useful:
\[
V_{oc} = \frac{E_{\text{CB}}}{q} - \frac{E_{\text{redox}}}{q} + \frac{kT}{q}\ln\left(\frac{n_{\text{CB}}}{N_{\text{CB}}}\right)
\]
\begin{itemize}
\item $n_{\text{CB}}$: electron density in TiO$_2$ at open circuit.
\item $N_{\text{CB}}$: effective density of states.
\end{itemize}
\textbf{Increasing $V_{oc}$:}
\begin{itemize}
\item Raise TiO$_2$ CB (e.g., MgO coating).
\item Use redox couples with more positive potential (Co$^{\text{III/II}}$).
\end{itemize}
\cite{Robertson2006}
\end{frame}


\begin{frame}{Fill Factor ($FF$)}
\frametitle{Fill Factor ($FF$)}

\textbf{Fill Factor ($FF$)} — Series and Shunt Resistance
\[
FF = \frac{P_{\max}}{J_{sc} V_{oc}}
\]
\textbf{Limiting factors:}
\begin{itemize}
\item \textbf{Series resistance ($R_s$):} FTO resistance, counter electrode catalysis, electrolyte conductivity. High $R_s$ reduces FF.
\item \textbf{Shunt resistance ($R_{sh}$):} Pinholes, direct contact between FTO and electrolyte. Low $R_{sh}$ reduces FF and $V_{oc}$.
\end{itemize}
Typical good DSSC: $R_s < 10\ \Omega\cdot$cm$^2$, $R_{sh} > 1000\ \Omega\cdot$cm$^2$ → $FF \approx 0.65–0.75$.
\cite{Snaith2007}
\end{frame}


\begin{frame}{Inorganic Sensitizers}
\frametitle{Inorganic Sensitizers}

\textbf{Inorganic Sensitizers} — Polypyridyl Ruthenium Complexes
\begin{itemize}
\item \textbf{N3:} [Ru(dcbpy)$_2$(NCS)$_2$] – 10.4\% efficiency.
\item \textbf{N719:} Di-tetrabutylammonium salt of N3 – 11.2\%.
\item \textbf{Black dye:} [Ru(tcterpy)(NCS)$_3$] – panchromatic absorption to 920 nm – 11.4\%.
\end{itemize}
\textbf{Pros:} Broad absorption, excellent stability, high efficiency.
\textbf{Cons:} Expensive, low extinction coefficient, contains heavy metal.
\cite{Nazeeruddin2005}
\end{frame}


\begin{frame}{Metal-Free Organic Dyes}
\frametitle{Metal-Free Organic Dyes}

\textbf{Metal-Free Organic Dyes} — D-$\pi$-A Architecture
\begin{center}
Donor (D) – $\pi$-bridge – Acceptor (A) + anchoring group
\end{center}
\textbf{Examples:}
\begin{itemize}
\item Coumarin derivatives (NKX-2677): 8–9\%.
\item Indoline (D149): high efficiency with Co-electrolyte.
\item Porphyrin (YD2-o-C8): 13\% record (2014).
\end{itemize}
\textbf{Advantages:}
\begin{itemize}
\item High $\varepsilon$ ($>5\times10^4$ M$^{-1}$cm$^{-1}$).
\item Tunable HOMO/LUMO.
\item Low cost, no noble metals.
\end{itemize}
\cite{Mishra2009}
\end{frame}


\begin{frame}{Natural Pigments}
\frametitle{Natural Pigments}

\textbf{Natural Pigments} — Anthocyanins and Chlorophyll
\begin{itemize}
\item \textbf{Anthocyanins} (berries, red cabbage): $\lambda_{\max} \approx 530$ nm, efficiency 0.5–1.5\%.
\item \textbf{Chlorophyll derivatives:} Better red absorption, but poor stability.
\end{itemize}
\textbf{Pros:} Biocompatible, extremely cheap, educational.
\textbf{Cons:} Narrow absorption, low efficiency, fast degradation.
\vspace{2mm}
\textit{Research use only} – not practical for commercial cells.
\cite{Calogero2012}
\end{frame}


\begin{frame}{Molar Extinction Coefficients}
\frametitle{Molar Extinction Coefficients}

\textbf{Molar Extinction Coefficients} — Why Organic Dyes Excel in Thin Layers
\begin{itemize}
\item Ru N719: $\varepsilon_{530} = 1.4 \times 10^4$ M$^{-1}$cm$^{-1}$.
\item Organic D149: $\varepsilon_{530} = 6.8 \times 10^4$ M$^{-1}$cm$^{-1}$.
\end{itemize}
\textbf{Consequence for film thickness:}
\[
\text{Absorbance} = \varepsilon \cdot c_{\text{dye}} \cdot L
\]
With high $\varepsilon$, a 2–4 \textmu m TiO$_2$ film suffices → lower recombination, faster transport.
\vspace{2mm}
Enables flexible, transparent DSSCs.
\cite{Grätzel2009}
\end{frame}


\begin{frame}{Anchoring Groups}
\frametitle{Anchoring Groups}

\textbf{Anchoring Groups} — Binding to TiO$_2$
\textbf{Most common:} Carboxylic acid (–COOH)
\begin{itemize}
\item Forms ester-like bond with Ti(IV) sites.
\item Provides electronic coupling for injection.
\end{itemize}
\textbf{Alternatives:}
\begin{itemize}
\item Phosphonic acid (–PO$_3$H$_2$): More stable, slower injection.
\item Siloxane: Stable but poor electronic coupling.
\end{itemize}
\textbf{Desorption risk:} In water or high pH, carboxylic anchors detach → stability issue.
\cite{Galoppini2004}
\end{frame}


\begin{frame}{Liquid Electrolytes}
\frametitle{Liquid Electrolytes}

\textbf{Liquid Electrolytes} — Acetonitrile and $I^-/I_3^-$
\textbf{Composition:}
\begin{itemize}
\item Solvent: acetonitrile (low viscosity) or valeronitrile (less volatile).
\item 0.5 M LiI, 0.05 M I$_2$, 0.5 M tert-butylpyridine (TBP) – improves $V_{oc}$.
\end{itemize}
\textbf{Advantages:} High conductivity ($\sim$10 mS/cm), good pore penetration.
\textbf{Drawbacks:} Volatile, leakage, solvent evaporation over time.
\vspace{2mm}
Record efficiency with liquid: 13\% (co-sensitized, 2022).
\cite{Yella2011}
\end{frame}


\begin{frame}{Alternative Redox Couples}
\frametitle{Alternative Redox Couples}

\textbf{Alternative Redox Couples} — Cobalt and Copper Complexes
\begin{itemize}
\item \textbf{Cobalt(III/II) tris(bipyridine):} $E^0 \approx +0.56$ V vs. NHE → $V_{oc}$ up to 1.0 V.
\item \textbf{Problem:} Slower dye regeneration; requires high extinction organic dyes.
\item \textbf{Copper(I/II) complexes:} Very fast regeneration, low recombination – promising $>$14\%.
\end{itemize}
\textbf{Recent record:} Cu-based electrolyte + organic dye → 14.3\% (2021, lab scale).
\cite{Saygili2018}
\end{frame}


\begin{frame}{Solid-State DSSCs}
\frametitle{Solid-State DSSCs}

\textbf{Solid-State DSSCs} — Spiro-OMeTAD as HTM
\begin{itemize}
\item Replace liquid electrolyte with organic hole transporter (e.g., Spiro-OMeTAD).
\item No leakage, easier encapsulation.
\item Challenges: Poor pore filling, faster recombination.
\end{itemize}
\textbf{Improvements:}
\begin{itemize}
\item Add LiTFSI and Co(tbp)$_2$(TFSI)$_2$ to dope Spiro-OMeTAD.
\item Record ssDSSC: 8.5\% (small area).
\end{itemize}
Still far behind liquid cells due to incomplete contact.
\cite{Bach1998}
\end{frame}


\begin{frame}{Quasi-Solid State}
\frametitle{Quasi-Solid State}

\textbf{Quasi-Solid State} — Polymer Gels and Ionic Liquids
\textbf{Approaches:}
\begin{itemize}
\item \textbf{Polymer gel electrolytes:} PMMA, PEO + liquid electrolyte → high viscosity.
\item \textbf{Ionic liquids:} 1-ethyl-3-methylimidazolium iodide – negligible vapor pressure.
\end{itemize}
\textbf{Performance:}
\begin{itemize}
\item Gel: 7–9\% efficiency, improved stability (1000h at 60°C).
\item Ionic liquid: 8–10\%, slower kinetics.
\end{itemize}
Best compromise for commercial flexible DSSCs.
\cite{Wang2003}
\end{frame}


\begin{frame}{Mechanism of Triiodide Reduction}
\frametitle{Mechanism of Triiodide Reduction}

\textbf{Mechanism of Triiodide Reduction} — At the Counter Electrode
Overall: $I_3^- + 2e^- \rightarrow 3I^-$
\vspace{2mm}
\textbf{Electrochemical steps (on Pt):}
\begin{enumerate}
\item $I_3^- + e^- \rightarrow I_2^- + I^-$ (rate-determining)
\item $I_2^- + e^- \rightarrow 2I^-$
\end{enumerate}
\textbf{Catalyst quality:} Measured by $R_{ct}$ from EIS (Electrochemical Impedance Spectroscopy).
\vspace{2mm}
Good Pt: $R_{ct} \approx 1\ \Omega$; poor carbon: $R_{ct} > 50\ \Omega$.
\cite{Hauch2001}
\end{frame}


\begin{frame}{Platinum (Pt)}
\frametitle{Platinum (Pt)}

\textbf{Platinum (Pt)} — The Gold Standard
\textbf{Preparation:}
\begin{itemize}
\item Sputtering: $<$10 nm, high cost.
\item Thermal decomposition: H$_2$PtCl$_6$ in isopropanol, 400°C – industrial.
\end{itemize}
\textbf{Performance:}
\begin{itemize}
\item Low $R_{ct}$, high exchange current density.
\item Stable in $I^-/I_3^-$.
\end{itemize}
\textbf{Drawback:} Pt dissolves in $I_3^-$ over long term (thousands of hours) – requires excess $I^-$.
\cite{Olsen2000}
\end{frame}


\begin{frame}{Carbon-Based Alternatives}
\frametitle{Carbon-Based Alternatives}

\textbf{Carbon-Based Alternatives} — Graphene, CNTs, Carbon Black
\textbf{Advantages:}
\begin{itemize}
\item Cheap, abundant, corrosion-resistant.
\item High surface area (active sites).
\end{itemize}
\textbf{Types:}
\begin{itemize}
\item Carbon black + graphite: $R_{ct} \approx 10–20\ \Omega$ – acceptable for low-cost cells.
\item Graphene: Superior conductivity, $R_{ct} \approx 5\ \Omega$.
\item CNT networks: Excellent 3D catalytic surface.
\end{itemize}
Efficiency comparable to Pt (10–11\%) with optimized carbon.
\cite{Kavan2011}
\end{frame}


\begin{frame}{Conducting Polymers (PEDOT:PSS)}
\frametitle{Conducting Polymers (PEDOT:PSS)}

\textbf{Conducting Polymers (PEDOT:PSS)} — Affordable and Flexible
\begin{itemize}
\item Poly(3,4-ethylenedioxythiophene):polystyrene sulfonate.
\item Solution-processable, transparent, flexible.
\item Catalytic for $I_3^-$ reduction (redox-active polymer).
\end{itemize}
\textbf{Performance:}
\begin{itemize}
\item $R_{ct} \approx 10–30\ \Omega$ (needs optimization).
\item Efficiency: 7–9\% in lab-scale.
\end{itemize}
Ideal for roll-to-roll printing on plastic.
\cite{Yoo2012}
\end{frame}


\begin{frame}{Power Conversion Efficiency (PCE)}
\frametitle{Power Conversion Efficiency (PCE)}

\textbf{Power Conversion Efficiency (PCE)} — Formula and Components
\[
\eta = \frac{J_{sc} \cdot V_{oc} \cdot FF}{P_{\text{in}}}
\]
where $P_{\text{in}} = 100\ \text{mW/cm}^2$ (AM 1.5G).
\vspace{2mm}
\textbf{Typical values for DSSC:}
\begin{itemize}
\item $J_{sc} = 15–22$ mA/cm$^2$
\item $V_{oc} = 0.7–0.9$ V
\item $FF = 0.65–0.75$
\item $\eta = 10–13\%$ (lab record)
\end{itemize}
\textbf{Shockley-Queisser limit for single-junction DSSC:} $\sim$32\%.
\cite{Hagfeldt2010}
\end{frame}


\begin{frame}{IPCE}
\frametitle{IPCE}

\textbf{IPCE} — Spectral Response Analysis
\[
\text{IPCE}(\lambda) = \frac{1240 \cdot J_{sc}(\lambda)}{\lambda \cdot P_{\text{in}}(\lambda)} \times 100\%
\]
\begin{itemize}
\item Measures how efficiently each wavelength is converted.
\item For DSSCs: IPCE = LHE($\lambda$) $\cdot$ $\phi_{inj}$ $\cdot$ $\phi_{coll}$.
\end{itemize}
\textbf{Typical IPCE:} 80–90\% at peak wavelength (550 nm) for Ru-dye cells.
\vspace{2mm}
Integration of IPCE over AM1.5G gives $J_{sc}$ – diagnostic tool.
\cite{Snaith2005}
\end{frame}


\begin{frame}{Light Harvesting Efficiency (LHE) and Quantum Yield}
\frametitle{Light Harvesting Efficiency (LHE) and Quantum Yield}

\textbf{Light Harvesting Efficiency (LHE) and Quantum Yield}
\[
\text{LHE}(\lambda) = 1 - 10^{-A(\lambda)} = 1 - 10^{-\varepsilon c L}
\]
\begin{itemize}
\item $A(\lambda)$: absorbance of dye-loaded film.
\end{itemize}
\textbf{Internal Quantum Efficiency (IQE)} = $\phi_{inj} \cdot \phi_{coll}$.
\begin{itemize}
\item Near 1.0 in good DSSCs (100\% of injected electrons collected).
\end{itemize}
\textbf{External Quantum Efficiency (EQE)} = IPCE = LHE $\cdot$ IQE.
\cite{Hagfeldt2008}
\end{frame}


\begin{frame}{Current Efficiency Records}
\frametitle{Current Efficiency Records}

\textbf{Current Efficiency Records} — Lab-Scale vs. Module
\begin{table}
\centering
\begin{tabular}{lll}
\toprule
Type & Efficiency & Reference \\
\midrule
Ru-dye (N719) + liquid & 11.2\% & Nazeeruddin (2005) \\
Black dye + liquid & 11.4\% & Nazeeruddin (2005) \\
Porphyrin (YD2-o-C8) + Co & 13.0\% & Mathew (2014) \\
Co-sensitized (organic + porphyrin) & 14.3\% & Kakiage (2015) \\
Module (substrate 10 cm$^2$) & 9.9\% & Sharp (2017) \\
\bottomrule
\end{tabular}
\end{table}
\textit{Note:} Perovskite cells have overtaken DSSC (25.7\%), but DSSC excels in stability and indoor light.
\cite{Green2022}
\end{frame}


\begin{frame}{Co-sensitization}
\frametitle{Co-sensitization}

\textbf{Co-sensitization} — Covering the Full Solar Spectrum
\begin{itemize}
\item Use two or more dyes with complementary absorption.
\item Example: Porphyrin (red/NIR) + organic D-$\pi$-A (blue/green).
\item Requires careful control to avoid energy transfer quenching.
\end{itemize}
\textbf{Record (2015, Kakiage):} Co-sensitized with alkoxy-substituted organic dyes + porphyrin → $J_{sc} = 19.5$ mA/cm$^2$, $V_{oc}=0.94$ V, $\eta=14.3\%$.
\cite{Kakiage2015}
\end{frame}


\begin{frame}{Light Scattering Layers}
\frametitle{Light Scattering Layers}

\textbf{Light Scattering Layers} — Trapping Photons
\begin{itemize}
\item Add a 4–5 \textmu m top layer of larger TiO$_2$ particles (200–400 nm).
\item Increases optical path length without more dye.
\item Enhances red/NIR absorption.
\end{itemize}
\textbf{Effect:} $J_{sc}$ increase by 10–30\%.
\vspace{2mm}
\textit{Drawback:} Thicker film → more recombination. Optimized cells use bilayer: 10 \textmu m transparent + 5 \textmu m scattering.
\cite{Hore2005}
\end{frame}


\begin{frame}{Plasmonic Enhancement}
\frametitle{Plasmonic Enhancement}

\textbf{Plasmonic Enhancement} — Metal Nanoparticles (Ag, Au)
\begin{itemize}
\item Localized surface plasmon resonance (LSPR) in Ag or Au nanoparticles (5–30 nm) embedded near dye.
\item Enhanced electromagnetic field → higher dye absorption.
\item Can also increase electron injection via hot-electron transfer.
\end{itemize}
\textbf{Results:} Up to 20\% increase in $J_{sc}$ reported.
\vspace{2mm}
\textit{Caution:} Metal can also act as recombination center if not insulated (e.g., SiO$_2$ shell).
\cite{Standridge2009}
\end{frame}


\begin{frame}{Surface Passivation}
\frametitle{Surface Passivation}

\textbf{Surface Passivation} — Atomic Layer Deposition (ALD) to Reduce Recombination
\begin{itemize}
\item TiO$_2$ surface has oxygen vacancies and defects → recombination sites for electrons with $I_3^-$.
\item Coating with ultrathin insulator (Al$_2$O$_3$, Nb$_2$O$_5$, MgO) reduces back-electron transfer.
\end{itemize}
\textbf{ALD Al$_2$O$_3$ (1–2 monolayers):}
\begin{itemize}
\item Increases $V_{oc}$ by 50–100 mV.
\item Preserves injection efficiency.
\end{itemize}
Standard technique in high-efficiency cells.
\cite{Palomares2004}
\end{frame}


\begin{frame}{Stability Challenges}
\frametitle{Stability Challenges}

\textbf{Stability Challenges} — Dye Desorption, Electrolyte Leakage, UV Degradation
\begin{enumerate}
\item \textbf{Dye desorption:} Polar solvents (water) displace anchoring group.
\item \textbf{Electrolyte leakage:} Liquid cells lose solvent → performance drop.
\item \textbf{UV degradation:} TiO$_2$ photocatalysis degrades dye; UV-filtering cover glass helps.
\end{enumerate}
\textbf{Test protocols:} ISOS-D-1 (dark, 85°C), ISOS-L-1 (light soaking at 60°C).
\vspace{2mm}
Best cells: 1000h at 80°C with $<$10\% loss.
\cite{Kato2009}
\end{frame}


\begin{frame}{Encapsulation Technologies}
\frametitle{Encapsulation Technologies}

\textbf{Encapsulation Technologies} — Glass-to-Glass and Thermoplastic Sealing
\textbf{Materials:}
\begin{itemize}
\item Surlyn (ionomer resin) – heat-sealed at 120°C.
\item UV-curable epoxy – for rigid substrates.
\end{itemize}
\textbf{Architecture:}
\begin{itemize}
\item Glass-FTO – TiO$_2$/dye – electrolyte – Pt/FTO – Glass.
\item Laser-welded glass frit for hermetic seal.
\end{itemize}
Flexible cells require barrier films (e.g., PET/Al multilayer) – still challenging.
\cite{Sastrawan2006}
\end{frame}


\begin{frame}{Unique Advantages}
\frametitle{Unique Advantages}

\textbf{Unique Advantages} — Low-Light and Indoor Performance
\begin{itemize}
\item DSSCs work efficiently under diffuse, low-intensity light (100–500 lux).
\item Indoor light (LED, fluorescent): 10–25\% efficiency under 200 lux (vs. Si 8–12\%).
\item Excellent for IoT sensors, wearables, indoor electronics.
\end{itemize}
\textbf{Why?} Low series resistance losses, no need for full sunlight to overcome diode ideality factor.
\vspace{2mm}
Commercial examples: Ambient Photonics (DSSC-powered keyboards, price tags).
\cite{Freitag2017}
\end{frame}


\begin{frame}{Building Integrated PV (BIPV)}
\frametitle{Building Integrated PV (BIPV)}

\textbf{Building Integrated PV (BIPV)} — Transparency and Aesthetic Color
\begin{itemize}
\item DSSCs can be made semitransparent (replace dye with scattering particles or thin film).
\item Colors tuned by dye choice (red, blue, green, purple).
\item Suitable for windows, facades, noise barriers.
\end{itemize}
\textbf{Efficiency-transparency trade-off:}
\begin{itemize}
\item Opaque: 10–12\%.
\item 30\% transparent: 6–8\%.
\end{itemize}
Pilot installations: Fraunhofer ISE (Germany), Sony (Japan).
\cite{Yoon2011}
\end{frame}


\begin{frame}{Summary and Outlook}
\frametitle{Summary and Outlook}

\textbf{Summary and Outlook} — Bridging Lab Efficiency and Commercial Durability
\textbf{Strengths:}
\begin{itemize}
\item Low-cost materials, non-vacuum processing.
\item Works indoors, flexible, colorful.
\end{itemize}
\textbf{Weaknesses:}
\begin{itemize}
\item Liquid electrolyte stability.
\item Efficiency lagging behind perovskites and Si.
\end{itemize}
\textbf{Future directions:}
\begin{itemize}
\item Solid-state DSSCs with robust HTMs.
\item Machine learning for dye/electrolyte discovery.
\item Niche markets: indoor IoT, BIPV, portable charging.
\end{itemize}
DSSCs will remain relevant as a sustainable, design-flexible PV technology.
\cite{Grätzel2021}
\end{frame}


\begin{frame}[allowframebreaks]{References (Page 1/4)}
\frametitle{References (Page 1/4)}
\vspace{0.5em}\begin{center}\textcolor{gray}{\small Page 1/4} \hfill \hyperlink{page.2}{\textcolor{blue}{\small Next Page →}}\end{center}

\begin{thebibliography}{99}
\bibitem{ORegan1991} B. O'Regan, M. Grätzel, \textit{Nature} \textbf{1991}, 353, 737. \hfill\textcolor{gray}{\tiny [Cited on slides: \hyperlink{page.3}{3}]}
\bibitem{Grätzel2001} M. Grätzel, \textit{Nature} \textbf{2001}, 414, 338.
\bibitem{Hagfeldt2010} A. Hagfeldt et al., \textit{Chem. Rev.} \textbf{2010}, 110, 6595. \hfill\textcolor{gray}{\tiny [Cited on slides: \hyperlink{page.7}{7}, \hyperlink{page.37}{37}]}
\bibitem{Green2022} M. Green et al., \textit{Prog. Photovolt.} \textbf{2022}, 30, 3. \hfill\textcolor{gray}{\tiny [Cited on slides: \hyperlink{page.2}{2}, \hyperlink{page.40}{40}]}
\bibitem{Peter2007} L. M. Peter, \textit{Phys. Chem. Chem. Phys.} \textbf{2007}, 9, 2630. \hfill\textcolor{gray}{\tiny [Cited on slides: \hyperlink{page.4}{4}]}
\bibitem{Kalyanasundaram2010} K. Kalyanasundaram, \textit{Energy Environ. Sci.} \textbf{2010}, 3, 919. \hfill\textcolor{gray}{\tiny [Cited on slides: \hyperlink{page.5}{5}]}
\bibitem{Gordillo2006} G. Gordillo et al., \textit{Sol. Energy Mater. Sol. Cells} \textbf{2006}, 90, 2556. \hfill\textcolor{gray}{\tiny [Cited on slides: \hyperlink{page.6}{6}]}
\bibitem{Grätzel2005} M. Grätzel, \textit{Inorg. Chem.} \textbf{2005}, 44, 6841. \hfill\textcolor{gray}{\tiny [Cited on slides: \hyperlink{page.8}{8}]}
\bibitem{Mathew2014} S. Mathew et al., \textit{Nat. Chem.} \textbf{2014}, 6, 242. \hfill\textcolor{gray}{\tiny [Cited on slides: \hyperlink{page.9}{9}]}
\bibitem{Feldt2010} S. M. Feldt et al., \textit{J. Am. Chem. Soc.} \textbf{2010}, 132, 16714. \hfill\textcolor{gray}{\tiny [Cited on slides: \hyperlink{page.10}{10}]}
\bibitem{Trancik2007} J. E. Trancik et al., \textit{Adv. Mater.} \textbf{2007}, 19, 657. \hfill\textcolor{gray}{\tiny [Cited on slides: \hyperlink{page.11}{11}]}
\bibitem{Ardo2009} S. Ardo, G. J. Meyer, \textit{Chem. Soc. Rev.} \textbf{2009}, 38, 115. \hfill\textcolor{gray}{\tiny [Cited on slides: \hyperlink{page.13}{13}]}
\end{thebibliography}
\end{frame}

\begin{frame}[allowframebreaks]{References (Page 2/4)}
\frametitle{References (Page 2/4)}
\vspace{0.5em}\begin{center}\hyperlink{page.1}{\textcolor{blue}{\small ← Previous Page}} \hfill \textcolor{gray}{\small Page 2/4} \hfill \hyperlink{page.3}{\textcolor{blue}{\small Next Page →}}\end{center}

\begin{thebibliography}{99}
\bibitem{Koops2009} S. E. Koops et al., \textit{J. Am. Chem. Soc.} \textbf{2009}, 131, 5218. \hfill\textcolor{gray}{\tiny [Cited on slides: \hyperlink{page.14}{14}]}
\bibitem{Nelson2003} J. Nelson, \textit{Phys. Rev. B} \textbf{2003}, 67, 155209. \hfill\textcolor{gray}{\tiny [Cited on slides: \hyperlink{page.15}{15}]}
\bibitem{Clifford2011} J. N. Clifford et al., \textit{J. Phys. Chem. C} \textbf{2011}, 115, 21847. \hfill\textcolor{gray}{\tiny [Cited on slides: \hyperlink{page.16}{16}]}
\bibitem{Hilgendorff2000} M. Hilgendorff, M. Grätzel, \textit{Electrochim. Acta} \textbf{2000}, 45, 4577. \hfill\textcolor{gray}{\tiny [Cited on slides: \hyperlink{page.17}{17}]}
\bibitem{Durrant2004} J. R. Durrant, \textit{J. Photochem. Photobiol. A} \textbf{2004}, 164, 23. \hfill\textcolor{gray}{\tiny [Cited on slides: \hyperlink{page.18}{18}]}
\bibitem{Frank2004} A. J. Frank et al., \textit{Coord. Chem. Rev.} \textbf{2004}, 248, 1165. \hfill\textcolor{gray}{\tiny [Cited on slides: \hyperlink{page.19}{19}]}
\bibitem{Cahen2003} D. Cahen et al., \textit{J. Phys. Chem. B} \textbf{2003}, 107, 10253. \hfill\textcolor{gray}{\tiny [Cited on slides: \hyperlink{page.20}{20}]}
\bibitem{Nazeeruddin2003} M. K. Nazeeruddin et al., \textit{J. Am. Chem. Soc.} \textbf{2003}, 125, 14522. \hfill\textcolor{gray}{\tiny [Cited on slides: \hyperlink{page.21}{21}]}
\bibitem{Robertson2006} N. Robertson, \textit{Angew. Chem. Int. Ed.} \textbf{2006}, 45, 2338. \hfill\textcolor{gray}{\tiny [Cited on slides: \hyperlink{page.22}{22}]}
\bibitem{Snaith2007} H. J. Snaith, M. Grätzel, \textit{Appl. Phys. Lett.} \textbf{2007}, 91, 092103. \hfill\textcolor{gray}{\tiny [Cited on slides: \hyperlink{page.23}{23}]}
\bibitem{Nazeeruddin2005} M. K. Nazeeruddin et al., \textit{J. Am. Chem. Soc.} \textbf{2005}, 127, 16835. \hfill\textcolor{gray}{\tiny [Cited on slides: \hyperlink{page.24}{24}]}
\bibitem{Mishra2009} A. Mishra et al., \textit{Angew. Chem. Int. Ed.} \textbf{2009}, 48, 2474. \hfill\textcolor{gray}{\tiny [Cited on slides: \hyperlink{page.25}{25}]}
\end{thebibliography}
\end{frame}

\begin{frame}[allowframebreaks]{References (Page 3/4)}
\frametitle{References (Page 3/4)}
\vspace{0.5em}\begin{center}\hyperlink{page.2}{\textcolor{blue}{\small ← Previous Page}} \hfill \textcolor{gray}{\small Page 3/4} \hfill \hyperlink{page.4}{\textcolor{blue}{\small Next Page →}}\end{center}

\begin{thebibliography}{99}
\bibitem{Calogero2012} G. Calogero et al., \textit{Energy Environ. Sci.} \textbf{2012}, 5, 9882. \hfill\textcolor{gray}{\tiny [Cited on slides: \hyperlink{page.26}{26}]}
\bibitem{Grätzel2009} M. Grätzel, \textit{Acc. Chem. Res.} \textbf{2009}, 42, 1788. \hfill\textcolor{gray}{\tiny [Cited on slides: \hyperlink{page.27}{27}]}
\bibitem{Galoppini2004} E. Galoppini, \textit{Coord. Chem. Rev.} \textbf{2004}, 248, 1283. \hfill\textcolor{gray}{\tiny [Cited on slides: \hyperlink{page.28}{28}]}
\bibitem{Yella2011} A. Yella et al., \textit{Science} \textbf{2011}, 334, 629. \hfill\textcolor{gray}{\tiny [Cited on slides: \hyperlink{page.29}{29}]}
\bibitem{Saygili2018} Y. Saygili et al., \textit{J. Mater. Chem. A} \textbf{2018}, 6, 15433. \hfill\textcolor{gray}{\tiny [Cited on slides: \hyperlink{page.30}{30}]}
\bibitem{Bach1998} U. Bach et al., \textit{Nature} \textbf{1998}, 395, 583. \hfill\textcolor{gray}{\tiny [Cited on slides: \hyperlink{page.31}{31}]}
\bibitem{Wang2003} P. Wang et al., \textit{J. Am. Chem. Soc.} \textbf{2003}, 125, 1166. \hfill\textcolor{gray}{\tiny [Cited on slides: \hyperlink{page.32}{32}]}
\bibitem{Hauch2001} A. Hauch, A. Georg, \textit{Electrochim. Acta} \textbf{2001}, 46, 3457. \hfill\textcolor{gray}{\tiny [Cited on slides: \hyperlink{page.33}{33}]}
\bibitem{Olsen2000} E. Olsen et al., \textit{Sol. Energy Mater. Sol. Cells} \textbf{2000}, 63, 267. \hfill\textcolor{gray}{\tiny [Cited on slides: \hyperlink{page.34}{34}]}
\bibitem{Kavan2011} L. Kavan, J. H. Yum, M. Grätzel, \textit{ACS Nano} \textbf{2011}, 5, 165. \hfill\textcolor{gray}{\tiny [Cited on slides: \hyperlink{page.35}{35}]}
\bibitem{Yoo2012} D. Yoo et al., \textit{Nano Energy} \textbf{2012}, 1, 414. \hfill\textcolor{gray}{\tiny [Cited on slides: \hyperlink{page.36}{36}]}
\bibitem{Snaith2005} H. J. Snaith, L. Schmidt-Mende, \textit{Adv. Mater.} \textbf{2005}, 17, 2851. \hfill\textcolor{gray}{\tiny [Cited on slides: \hyperlink{page.38}{38}]}
\end{thebibliography}
\end{frame}

\begin{frame}[allowframebreaks]{References (Page 4/4)}
\frametitle{References (Page 4/4)}
\vspace{0.5em}\begin{center}\hyperlink{page.3}{\textcolor{blue}{\small ← Previous Page}} \hfill \textcolor{gray}{\small Page 4/4}\end{center}

\begin{thebibliography}{99}
\bibitem{Hagfeldt2008} A. Hagfeldt, M. Grätzel, \textit{Acc. Chem. Res.} \textbf{2008}, 33, 269. \hfill\textcolor{gray}{\tiny [Cited on slides: \hyperlink{page.39}{39}]}
\bibitem{Kakiage2015} K. Kakiage et al., \textit{Chem. Commun.} \textbf{2015}, 51, 6315. \hfill\textcolor{gray}{\tiny [Cited on slides: \hyperlink{page.41}{41}]}
\bibitem{Hore2005} S. Hore et al., \textit{Chem. Commun.} \textbf{2005}, 40, 511. \hfill\textcolor{gray}{\tiny [Cited on slides: \hyperlink{page.42}{42}]}
\bibitem{Standridge2009} S. D. Standridge et al., \textit{J. Am. Chem. Soc.} \textbf{2009}, 131, 18428. \hfill\textcolor{gray}{\tiny [Cited on slides: \hyperlink{page.43}{43}]}
\bibitem{Palomares2004} E. Palomares et al., \textit{J. Am. Chem. Soc.} \textbf{2004}, 126, 12250. \hfill\textcolor{gray}{\tiny [Cited on slides: \hyperlink{page.44}{44}]}
\bibitem{Kato2009} N. Kato et al., \textit{Sol. Energy Mater. Sol. Cells} \textbf{2009}, 93, 893. \hfill\textcolor{gray}{\tiny [Cited on slides: \hyperlink{page.45}{45}]}
\bibitem{Sastrawan2006} R. Sastrawan et al., \textit{Sol. Energy Mater. Sol. Cells} \textbf{2006}, 90, 1680. \hfill\textcolor{gray}{\tiny [Cited on slides: \hyperlink{page.46}{46}]}
\bibitem{Freitag2017} M. Freitag et al., \textit{Nat. Photonics} \textbf{2017}, 11, 372. \hfill\textcolor{gray}{\tiny [Cited on slides: \hyperlink{page.47}{47}]}
\bibitem{Yoon2011} J. H. Yoon et al., \textit{Energy Build.} \textbf{2011}, 43, 1735. \hfill\textcolor{gray}{\tiny [Cited on slides: \hyperlink{page.48}{48}]}
\bibitem{Grätzel2021} M. Grätzel, \textit{Acc. Chem. Res.} \textbf{2021}, 54, 3015. \hfill\textcolor{gray}{\tiny [Cited on slides: \hyperlink{page.49}{49}]}
\end{thebibliography}
\end{frame}

\end{document}
